公开研究情绪高度一致：AI capex 是 2026 年硬科技资产最强主线之一，市场更愿意给直接兑现的光模块、AI PCB、服务器 ODM 和液冷/供配电龙头估值溢价。本轮已从东方财富公开研报接口归档并抽取核心券商 PDF，并为 300476 额外归档同花顺 iFinD 公开一致预期快照。原 18 个目标价标的均有可审计 broker/Street 锚；其中 17 个标的来自原始 PDF，胜宏科技来自 iFinD 快照，目标价区间 360.00--403.42 元、均值 381.71 元，进入 capped 10\% broker/Street 锚。扩展模型新增 38 个可发布目标价/公允价值标的，其中 13 个有明示券商目标价，24 个使用 AStock 自建公允价值，1 个使用 PS/SOTP；另有 3 个因盈利或模型分母不足降级观察。


\begin{exhibitbox}[公开研究情绪归纳]
\begin{tabularx}{\textwidth}{L{3.0cm}X X}
\toprule
\textbf{共识方向} & \textbf{支持证据} & \textbf{本报告处理}\\
\midrule
AI capex 继续扩张 & NVIDIA、Dell'Oro、JLL 与国家数据局政策口径均支持需求强度 & 用作行业需求锚，不直接等同于单公司收入\\
光互联景气强 & LightCounting 指向 AI 集群光互联长期空间，但也提示供应链均衡风险 & 龙头给盈利信用，器件/扩散标的给可选性信用\\
液冷和供电升级 & JLL、NVIDIA 架构、英维克和科华公告共同指向高密机柜约束 & 只在产品/认证/项目证据明确时给估值权重\\
AIDC 运营商重估 & 润泽披露 AIDC 增长、200MW/液冷项目和客户结构优化 & 估值必须同时看上架率、现金流和折旧\\
\bottomrule
\end{tabularx}
\sourcenote{S01--S12；data/consensus\_analysis.md。}
\end{exhibitbox}

因此，本报告把“公开情绪”纳入综合目标价的市场锚，但权重低于基本面锚。若后续成交额和趋势继续强化而财务也兑现，市场锚权重可以上调；若订单、毛利或现金流低于预期，市场锚必须下调。
